% =============================================================
%  CHAPTER: Methodology
%  Thesis: Hallucination Detection in Large Language Models
%          via Modern Hopfield Retrieval over SwiGLU Memory Banks
%  Status: Working draft — final pipeline specification
% =============================================================

\chapter{Methodology}
\label{chap:methods}

This chapter specifies the methodology of the thesis. It translates the
theoretical framework of Chapter~\ref{chap:theory} into a concrete software
artifact, an operational pipeline, an evaluation protocol, and a final
experimental configuration that will be executed to test the central
hypothesis. The chapter is organised as a complete problem-solving cycle:
problem description, problem analysis, solution design, implementation,
testing, the proposed detection model, the statistical analysis method, and
the experimental setup of the final pipeline.

The presentation is intentionally constructive. Each phase introduces the
choices that shape the next phase, and the final section consolidates all of
those choices into a single executable plan. Earlier exploratory runs used
during the development of the framework are not discussed here — they served
only to dimension the problem; the configuration described below is the
definitive one and is the one that will produce the results reported in
Chapter~\ref{chap:experiments}.

% =============================================================
\section{Methodological Overview}
\label{sec:method_overview}
% =============================================================

The methodology follows the engineering cycle commonly used in
research-software projects: \emph{problem definition} $\rightarrow$
\emph{analysis} $\rightarrow$ \emph{design} $\rightarrow$
\emph{implementation} $\rightarrow$ \emph{testing} $\rightarrow$
\emph{evaluation}. Each step is reified as a section in this chapter and as
a concrete deliverable in the accompanying repository.

\begin{enumerate}[leftmargin=*]
    \item \textbf{Problem description}
        (Section~\ref{sec:problem_description}).
        Statement of the detection task, the binary target, the operating
        constraints, and the scope.
    \item \textbf{Problem analysis}
        (Section~\ref{sec:problem_analysis}).
        Decomposition of the task into measurable sub-problems: query
        identification, memory definition, retrieval scoring, divergence
        quantification, and label acquisition.
    \item \textbf{Solution design}
        (Section~\ref{sec:solution_design}).
        Formal definition of the proposed probe: the gated SwiGLU memory
        bank, the Modern Hopfield retrieval energy, the per-token
        information-theoretic divergences, and the inverse-temperature
        calibration.
    \item \textbf{Implementation}
        (Section~\ref{sec:implementation}).
        Description of the six-stage artefact-based pipeline and of the
        package modules that realise it.
    \item \textbf{Testing}
        (Section~\ref{sec:testing}).
        Unit, integration, and pipeline smoke tests; numerical
        regression tests on small fixed inputs.
    \item \textbf{Proposed model}
        (Section~\ref{sec:proposed_model}).
        The per-layer feature vector, the logistic-regression probe, and
        the held-out evaluation procedure.
    \item \textbf{Analysis method}
        (Section~\ref{sec:analysis_method}).
        Statistical procedures used to claim signal: nested
        cross-validation, bootstrap confidence intervals, permutation and
        within-category null distributions, scan correction across layers,
        paired-bootstrap nested-model tests, and effect-size reporting.
    \item \textbf{Experimental configuration}
        (Section~\ref{sec:experimental_config}).
        Datasets, models, splits, hyperparameters, baselines, and the
        ordered execution plan of the final pipeline.
\end{enumerate}

The pipeline is artefact-based: every stage writes a typed file to disk and
every downstream stage reads from that file. This makes the methodology
auditable, re-runnable in pieces, and resistant to silent regressions, as
the same artefacts are consumed by tests and by analysis notebooks.

% =============================================================
\section{Problem Description}
\label{sec:problem_description}
% =============================================================

% ----------------------------------------------------------------
\subsection{Detection task}
\label{subsec:task}
% ----------------------------------------------------------------

Let $f_\theta$ be a frozen, instruction-tuned, decoder-only language model
with $L$ layers and SwiGLU feed-forward sublayers
(Section~\ref{subsec:ffn}). Given a textual question $q$ and a set of gold
short-form answers $\mathcal{A}_q = \{a_1, \ldots, a_m\}$ produced by a
factoid question-answering dataset, $f_\theta$ generates an answer
$\hat{y} = f_\theta(q)$ by greedy autoregressive decoding under a fixed
chat-template wrapper.

A \emph{hallucination} is operationalised as a binary indicator $z_q \in
\{0, 1\}$ such that $z_q = 1$ when the generated answer $\hat{y}$ is not
factually supported by $\mathcal{A}_q$, and $z_q = 0$ otherwise. The label
$z_q$ is obtained by an external judging procedure described in
Section~\ref{subsec:labels}, not by the model itself.

The detection task is to construct a scoring function
\begin{equation}
    s : (q,\,\hat{y},\,\mathcal{I}(f_\theta, q, \hat{y})) \,\longmapsto\, \mathbb{R},
    \label{eq:detection_function}
\end{equation}
where $\mathcal{I}(\cdot)$ denotes a fixed family of internal observables
collected from the forward passes of $f_\theta$ on $q$ and during the
generation of $\hat{y}$. Higher values of $s$ indicate higher predicted
likelihood of $z_q = 1$. Quality is measured by the area under the receiver
operating characteristic curve (AUROC) of $s$ against $z_q$.

% ----------------------------------------------------------------
\subsection{Operating constraints}
\label{subsec:constraints}
% ----------------------------------------------------------------

The methodology adopts six constraints that scope what is allowed:

\begin{enumerate}[leftmargin=*]
    \item \emph{White-box access.} The model weights and activations are
          accessible. No fine-tuning of $f_\theta$ is permitted.
    \item \emph{Single forward pass for the main probe.} The Hopfield
          features are computed from one prefill plus one greedy
          generation; sampling-based baselines may use more.
    \item \emph{Frozen architecture.} The detection mechanism may not alter
          the decoding distribution. Decoded text under detection and
          under no detection must be identical.
    \item \emph{Mechanistic interpretability.} Each feature must admit a
          per-layer interpretation and must be traceable to a specific
          weight matrix or hidden state.
    \item \emph{Reproducibility.} The pipeline must be deterministic for
          fixed random seeds, fixed model checkpoints, and fixed
          chat-template configurations.
    \item \emph{Modest hardware envelope.} Stages requiring GPUs must be
          runnable on a single A100 40\,GB or, with 4-bit quantisation, on
          a consumer RTX-class card.
\end{enumerate}

These constraints align the methodology with the constraints of practical
deployment scenarios (local inference, no retraining, audit-friendly
features) and bound the engineering scope of the pipeline.

% ----------------------------------------------------------------
\subsection{Scope}
\label{subsec:scope}
% ----------------------------------------------------------------

The scope of this thesis is limited to (i) factoid question answering with
short gold answers, (ii) decoder-only models with SwiGLU MLP sublayers,
(iii) English-language datasets, and (iv) hallucinations that manifest in
the generated answer rather than in upstream chain-of-thought tokens. These
restrictions define the regime in which the proposed signal is claimed to
apply; cross-task generalisation is treated as out-of-scope.

% =============================================================
\section{Problem Analysis}
\label{sec:problem_analysis}
% =============================================================

The detection task decomposes into five well-defined sub-problems, each of
which corresponds to a specific module in the implementation.

% ----------------------------------------------------------------
\subsection{Query identification}
\label{subsec:query_id}
% ----------------------------------------------------------------

The methodology operates on the hidden state $\mathbf{h}_t^{(l)} \in
\mathbb{R}^D$ entering the SwiGLU feed-forward sublayer of layer $l$ at
position $t$ (the \emph{probe point} of
Figure~\ref{fig:transformer_block}). The first sub-problem is to specify
which token positions contribute the query during prefill and during
generation, and how to reduce a sequence of $T$ such states into the
quantities used downstream.

Two distinct positions matter:

\begin{itemize}[leftmargin=*]
    \item \textbf{Prefill reference.} The query is taken as the hidden
          state of the \emph{last prompt token} after the chat template has
          been applied. This is the state from which the model commits to
          producing $\hat{y}$, and it is a single token's representation
          rather than a mixture; this single-token choice is what makes
          the prefill reference comparable to the per-token generation
          states and avoids the mixture artefacts of mean-pooling across
          the prompt.
    \item \textbf{Generation queries.} During generation, the query at
          position $t$ is the hidden state $\mathbf{h}_t^{(l)}$ at the
          $t$-th newly generated token, captured once per token by a
          forward hook on the MLP module.
\end{itemize}

Both passes apply the same chat template, so the prefill last-token state
and the generation states share the same prompt prefix in their respective
attention contexts. This invariance is enforced in the implementation
(Section~\ref{sec:implementation}) and is verified by a smoke test.

% ----------------------------------------------------------------
\subsection{Memory definition}
\label{subsec:memory_def}
% ----------------------------------------------------------------

The second sub-problem is to specify the per-layer memory bank
$\mathbf{M}^{(l)} \in \mathbb{R}^{K_l \times D}$
(Equation~\ref{eq:memory_bank_definition}). The bank is a derived,
non-trainable function of the model parameters.

In SwiGLU, the activation coefficient of the $i$-th intermediate unit is
\begin{equation}
    a_i^{(l)}(\mathbf{h})
    =
    \mathrm{Swish}\!\bigl(\mathbf{h}\,{\mathbf{W}_1^{(l)}}^{\!\top}[i]\bigr)
    \,\cdot\,
    \bigl(\mathbf{h}\,{\mathbf{V}^{(l)}}^{\!\top}[i]\bigr).
    \label{eq:swiglu_unit_activation}
\end{equation}
Identifying the memory with either $\mathbf{W}_1^{(l)}$ or $\mathbf{V}^{(l)}$ alone
captures only half of the gating computation. The methodology therefore
defines the \emph{gated key bank}
\begin{equation}
    \mathbf{M}^{(l)} = \mathbf{W}_1^{(l)} \odot \mathbf{V}^{(l)}
    \in \mathbb{R}^{K_l \times D},
    \label{eq:gated_key_bank}
\end{equation}
where $\odot$ is the element-wise product. The $i$-th row
$\mathbf{m}_i^{(l)} = \mathbf{W}_1^{(l)}[i] \odot \mathbf{V}^{(l)}[i]$ is the
\emph{linearised gated key} of intermediate unit $i$, in the sense that
$\mathbf{m}_i^{(l)} \cdot \mathbf{h}$ approximates the gated activation
$a_i^{(l)}(\mathbf{h})$ around a small neighbourhood of the operating
point.\footnote{The exact $a_i^{(l)}(\mathbf{h})$ involves a non-linearity on
the gate branch and is therefore not a single inner product. The linearised
form is used as the retrieval score because the downstream Hopfield
machinery operates on dot products, and because the element-wise product
of the two projections is the unique parameter-derived quantity that mixes
both branches of the gate without changing the dimensionality.}

This bank is the principled key for SwiGLU. The implementation also
supports the un-gated bank
$\mathbf{M}^{(l)}_{\mathrm{up}} = \mathbf{V}^{(l)}$ as an ablation, and a
value-space bank
$\mathbf{M}^{(l)}_{\mathrm{val}} = (\mathbf{W}_2^{(l)})^{\!\top}$ for a
secondary probe at the FFN output.

Rows of the bank are L2-normalised after construction so that the
maximum row norm $M_l = \max_i \|\mathbf{m}_i^{(l)}\|_2$ equals $1$ for
every layer and the cross-layer constant $\tfrac{1}{2}M_l^2$ of
Equation~\eqref{eq:modern_hopfield_energy} is identically $\tfrac{1}{2}$.
This normalisation removes a layer-dependent additive offset from the
energy without affecting the relative ordering of layers within a sample,
and it makes the per-layer energies cross-layer comparable.

% ----------------------------------------------------------------
\subsection{Retrieval scoring}
\label{subsec:retrieval_scoring}
% ----------------------------------------------------------------

The third sub-problem is to compute the Modern Hopfield retrieval energy
$E(\mathbf{h}_t^{(l)};\,\mathbf{M}^{(l)},\,\beta)$ and the induced
retrieval distribution $\{p_i^{(l,t)}\}$
(Equations~\ref{eq:modern_hopfield_energy} and
\ref{eq:retrieval_distribution}). Two design points are fixed:

\begin{itemize}[leftmargin=*]
    \item \emph{Numerical stability.} The log-sum-exp term is computed by
          \texttt{torch.logsumexp}, not by naive exponentiation, so that
          the score range $[\beta \min_i s_i,\,\beta \max_i s_i]$ does not
          overflow at the operating $\beta$.
    \item \emph{Query mode.} The energy is computed in \texttt{dot} mode
          on the raw hidden state, with \texttt{normalize\_query=False}.
          This keeps the quadratic term $\tfrac{1}{2}\|\mathbf{h}_t^{(l)}\|^2$
          empirically varying and therefore makes the energy decomposition
          falsifiable; under unit-norm queries the quadratic term collapses
          to a constant and the empirical content of the LSE-vs-norm
          decomposition is lost.
\end{itemize}

The retrieval distribution is reported alongside the energy as the
softmax over the unbounded scores
$\beta\,\mathbf{m}_i^{(l)\!\top}\mathbf{h}_t^{(l)}$ and is the object on
which the divergences of Section~\ref{sec:divergences} act.

% ----------------------------------------------------------------
\subsection{Divergence quantification}
\label{subsec:divergence_quant}
% ----------------------------------------------------------------

The fourth sub-problem is to quantify how the per-layer retrieval
distribution shifts from the prefill reference to each generation token.
For every layer $l$ and every generation position $t$, the methodology
records the Jensen–Shannon divergence and the squared Hellinger distance
between the generation distribution $\{p_i^{(l,t)}\}$ and the prefill
reference distribution $\{p_i^{(l,\mathrm{ref})}\}$
(Equations~\ref{eq:js_divergence} and \ref{eq:hellinger_squared}). These
two divergences are symmetric, bounded, and have complementary
sensitivities: JS captures entropy-style mass redistribution and Hellinger
captures geometric distance between square-root profiles.

The Kullback–Leibler divergence is computed as the building block of JS
and reported as a diagnostic in the asymmetric form
$D_{\mathrm{KL}}(p^{(l,t)} \,\|\, p^{(l,\mathrm{ref})})$, but it is not used
as a feature.

% ----------------------------------------------------------------
\subsection{Label acquisition}
\label{subsec:labels}
% ----------------------------------------------------------------

The fifth sub-problem is to obtain reliable binary labels $z_q$ at scale.
Two complementary labelling procedures are used:

\begin{enumerate}[leftmargin=*]
    \item \textbf{Token-F1 heuristic.} The generated answer $\hat{y}$ and
          every gold answer $a_j$ are tokenised by the same tokenizer
          used by $f_\theta$. The maximum SQuAD-style token-F1 score
          over $\mathcal{A}_q$ is compared to a fixed threshold
          ($\tau_{\mathrm{F1}} = 0.3$). The heuristic provides an
          inexpensive, deterministic, but semantically blunt label.
    \item \textbf{LLM judge.} A separate language model is asked, via a
          locked prompt template, whether $\hat{y}$ is factually supported
          by $\mathcal{A}_q$. The judge's verbatim output is recorded
          alongside a parsed binary label and the SHA-1 of the prompt
          template.
\end{enumerate}

The two labels are reconciled by reporting Cohen's $\kappa$ as a
calibration statistic. The judge label is taken as the primary supervision
signal; the heuristic label is retained as a transparent cross-check and
as the supervision used in low-API-budget ablations. Both labels are
written as separate JSON artefacts and never overwrite each other, so the
choice of supervision can be changed without re-running the labelling
stage.

% =============================================================
\section{Solution Design}
\label{sec:solution_design}
% =============================================================

This section consolidates the sub-problems of Section~\ref{sec:problem_analysis}
into the formal definition of the proposed probe and into the algorithms
that produce its features.

% ----------------------------------------------------------------
\subsection{Per-layer feature design}
\label{subsec:feature_design}
% ----------------------------------------------------------------

For each layer $l \in \{0, \ldots, L-1\}$ the methodology defines a
small, fixed set of per-sample features. Let $\mathbf{h}^{(l,\mathrm{ref})}$
denote the prefill last-token hidden state, and let
$\{\mathbf{h}_t^{(l)}\}_{t=1}^{T_g}$ denote the generation-token hidden
states over the $T_g$ tokens produced for question $q$. Let
$E^{(l,t)}$, $H^{(l,t)}$, and $\bar{H}^{(l,t)}$ denote the energy, entropy,
and normalised entropy at $(l,t)$, and let
$\mathrm{JS}^{(l,t)}$ and $H^2_{\mathrm{Hel}}{}^{(l,t)}$ denote the JS and
squared-Hellinger divergences between $\{p_i^{(l,t)}\}$ and
$\{p_i^{(l,\mathrm{ref})}\}$.

For every layer $l$ the following scalars are computed:
\begin{align}
    \Delta E^{(l)}
        &= \tfrac{1}{T_g^\star} \!\sum_{t \in \mathcal{T}_q^\star}\! E^{(l,t)} - E^{(l,\mathrm{ref})}, \\
    \Delta H^{(l)}
        &= \tfrac{1}{T_g^\star} \!\sum_{t \in \mathcal{T}_q^\star}\! H^{(l,t)} - H^{(l,\mathrm{ref})}, \\
    \Delta \bar{H}^{(l)}
        &= \tfrac{1}{T_g^\star} \!\sum_{t \in \mathcal{T}_q^\star}\! \bar{H}^{(l,t)} - \bar{H}^{(l,\mathrm{ref})}, \\
    \overline{\mathrm{JS}}^{(l)}
        &= \tfrac{1}{T_g^\star} \!\sum_{t \in \mathcal{T}_q^\star}\! \mathrm{JS}^{(l,t)}, \\
    \overline{H^2_{\mathrm{Hel}}}{}^{(l)}
        &= \tfrac{1}{T_g^\star} \!\sum_{t \in \mathcal{T}_q^\star}\! H^2_{\mathrm{Hel}}{}^{(l,t)}.
    \label{eq:per_layer_features}
\end{align}
Here $\mathcal{T}_q^\star \subseteq \{1, \ldots, T_g\}$ is the
\emph{answer-span} subset of generated positions and $T_g^\star =
|\mathcal{T}_q^\star|$. The span is obtained by aligning $\hat{y}$ against
the gold answers $\mathcal{A}_q$ at the tokenizer level; when no
alignment is found, the span defaults to all content tokens (those that
remain after a deterministic punctuation/whitespace filter). This pooling
restriction is motivated by the observation that the detection signal of
interest is the model's commitment to factual claims, which lives in the
answer-bearing tokens rather than in connectors or punctuation.

The per-layer vector for sample $q$ is the concatenation
\begin{equation}
    \boldsymbol{\phi}_q
    =
    \bigl[
        \Delta E^{(0)}, \ldots, \Delta E^{(L-1)},\;
        \Delta \bar{H}^{(0)}, \ldots, \Delta \bar{H}^{(L-1)},\;
        \overline{\mathrm{JS}}^{(0)}, \ldots, \overline{\mathrm{JS}}^{(L-1)},\;
        \overline{H^2_{\mathrm{Hel}}}{}^{(0)}, \ldots, \overline{H^2_{\mathrm{Hel}}}{}^{(L-1)}
    \bigr]
    \in \mathbb{R}^{4L}.
    \label{eq:feature_vector}
\end{equation}
The ablation features $\Delta H^{(l)}$ (un-normalised entropy) are stored
as alternates and used only when an explicit ablation is reported.

% ----------------------------------------------------------------
\subsection{Inverse-temperature calibration}
\label{subsec:beta_calibration}
% ----------------------------------------------------------------

The retrieval sharpness $\beta$ controls whether the softmax over
$\mathbf{M}^{(l)}$ is near-uniform (no signal), maximally informative, or
collapsed onto a single row (no signal). Both extremes produce noisy
divergences. The methodology calibrates $\beta$ once per model–bank pair
by the following procedure:

\begin{enumerate}[leftmargin=*]
    \item Draw a fixed warm-up subset of $N_c = 30$ questions from the
          training split.
    \item Run prefill on each warm-up question and capture the raw
          scores $\mathbf{M}^{(l)}\mathbf{h}^{(l,\mathrm{ref})}$ at every
          layer.
    \item For each candidate $\beta \in \{1, 2, 5, 10, 15, 20, 30, 50,
          100\}$, evaluate the mean normalised entropy of
          $\mathrm{softmax}(\beta \cdot \mathrm{scores})$ across all
          captured layers and warm-up samples.
    \item Choose $\beta^\star$ as the value at which the mean normalised
          entropy crosses the operating target
          $\bar{H}^\star = 0.55$. Linear interpolation in
          $\log \beta$ is used between the two bracketing candidates.
\end{enumerate}

The choice $\bar{H}^\star = 0.55$ sits inside the \emph{detection band}:
above $\bar{H} \approx 0.7$ the softmax is nearly uniform and the
divergences vanish; below $\bar{H} \approx 0.3$ the softmax is collapsed
and the divergences saturate. This calibration is computed only on the
training split so that no test-set information leaks into $\beta^\star$.

% ----------------------------------------------------------------
\subsection{Algorithmic summary}
\label{subsec:algorithms}
% ----------------------------------------------------------------

The complete per-sample procedure is summarised in
Algorithm~\ref{alg:per_sample}. The procedure is purely functional in the
sense that it takes the model, the bank, and the question as inputs and
returns the feature vector $\boldsymbol{\phi}_q$ together with the
generated text $\hat{y}$; no in-place modifications of the model are
performed.

\begin{algorithm}[H]
\caption{Per-sample feature extraction}
\label{alg:per_sample}
\begin{algorithmic}[1]
\Require frozen model $f_\theta$, banks $\{\mathbf{M}^{(l)}\}_{l=0}^{L-1}$,
         inverse temperature $\beta^\star$, question $q$, gold answers $\mathcal{A}_q$
\State $p \gets \texttt{apply\_chat\_template}(q)$
\State Run prefill forward pass on $p$; at every layer $l$, capture the
       hidden state $\mathbf{h}^{(l,\mathrm{ref})}$ at the last prompt token
\State Compute $E^{(l,\mathrm{ref})}$, $H^{(l,\mathrm{ref})}$,
       $\bar{H}^{(l,\mathrm{ref})}$ and the reference distribution
       $\{p_i^{(l,\mathrm{ref})}\}$ for all $l$ via
       Equations~\eqref{eq:modern_hopfield_energy}--\eqref{eq:retrieval_distribution}
\State Greedy-decode $\hat{y}$ from $p$ with $T_g \le T_{\max}$ tokens; at
       every generated position $t$, capture
       $\{\mathbf{h}_t^{(l)}\}_{l=0}^{L-1}$
\For{each $(l, t)$}
    \State Compute $E^{(l,t)}$, $H^{(l,t)}$, $\bar{H}^{(l,t)}$,
           $\{p_i^{(l,t)}\}$
    \State Compute $\mathrm{JS}^{(l,t)}$ and
           $H^2_{\mathrm{Hel}}{}^{(l,t)}$ against
           $\{p_i^{(l,\mathrm{ref})}\}$
\EndFor
\State Align $\hat{y}$ to $\mathcal{A}_q$ at the tokenizer level to obtain
       $\mathcal{T}_q^\star$
\State Pool the per-token quantities over $\mathcal{T}_q^\star$ as in
       Equation~\eqref{eq:per_layer_features}
\State Assemble $\boldsymbol{\phi}_q$ as in Equation~\eqref{eq:feature_vector}
\State \Return $(\boldsymbol{\phi}_q,\,\hat{y})$
\end{algorithmic}
\end{algorithm}

% =============================================================
\section{Implementation}
\label{sec:implementation}
% =============================================================

The methodology is realised as a Python package, \texttt{hopfield\_llm},
organised around a six-stage artefact-based pipeline. Each stage produces a
typed artefact on disk and consumes only artefacts written by previous
stages. Stages 1–4 form the core capture-and-aggregate flow; Stages 5–6
are evaluation layers that may be re-run independently. The pipeline is
shown schematically in Figure~\ref{fig:pipeline_overview}.

\begin{figure}[H]
    \centering
    \begin{minipage}{0.92\linewidth}
    \texttt{%
    \begin{tabbing}
    \=[1]\quad\=\texttt{build-banks}\quad\quad\=GPU\quad\=model \=\(\rightarrow\) banks.pt\\
    [2]\>\texttt{run-trajectory}\>GPU\>banks + dataset\>\(\rightarrow\) \{id\}.npz, \{id\}.json\\
    [3]\>\texttt{analyze}\>CPU\>trajectory dir\>\(\rightarrow\) analysis.json\\
    [4]\>\texttt{visualize}\>CPU\>analysis.json\>\(\rightarrow\) plots/*.png\\
    [5]\>\texttt{label}\>CPU\>trajectory dir\>\(\rightarrow\) \{id\}\_labeled.json, calibration.json\\
    [6]\>\texttt{evaluate}\>CPU\>trajectory + labels\>\(\rightarrow\) report.html
    \end{tabbing}
    }
    \end{minipage}
    \caption{The six-stage pipeline. GPU-bound stages run once per
    \emph{(model, dataset, probe)} triple; CPU-bound stages can be
    re-executed against the same trajectories without reusing the GPU.}
    \label{fig:pipeline_overview}
\end{figure}

% ----------------------------------------------------------------
\subsection{Stage 1 — Bank extraction}
\label{subsec:stage1}
% ----------------------------------------------------------------

\textbf{Module: \texttt{memory/banks.py}.} The stage walks every layer of
the loaded model, resolves the SwiGLU sub-modules through a model-alias
registry (\texttt{gate\_proj}/\texttt{up\_proj} for Qwen and Llama;
\texttt{w1}/\texttt{w3} aliases for Mistral-family naming), and constructs
the requested bank. For the gated key bank
(Equation~\ref{eq:gated_key_bank}), the dequantised \texttt{gate\_proj} and
\texttt{up\_proj} weights are loaded into \texttt{float32} on CPU, their
element-wise product is computed, and the resulting matrix
$\mathbf{W}_{\mathrm{gate}}^{(l)} \odot \mathbf{V}^{(l)}$ is L2-normalised
row-wise. The output artefact \texttt{banks.pt} contains
$\{l \mapsto \texttt{Tensor}[K_l, D]\}$ together with a companion
\texttt{banks\_metadata.json} that records per-layer
$\{K_l, D, M_l^{\mathrm{raw}}, M_l^{\mathrm{used}}, \text{source}\}$.

The supported bank values are \texttt{gated\_key} (primary, principled
SwiGLU key), \texttt{up} (un-gated key, kept for ablation), and
\texttt{down\_values} (rows of $\mathbf{W}_2^\top$ for the value-space
ablation). The raw \texttt{down\_proj} weight is explicitly rejected with
a hard error, because its rows live in the output dimension and would
silently confuse keys with values.

% ----------------------------------------------------------------
\subsection{Stage 2 — Trajectory capture}
\label{subsec:stage2}
% ----------------------------------------------------------------

\textbf{Modules: \texttt{hooks/capture.py}, \texttt{pipeline/stages.py}.}
For each sample, the stage runs two forward passes: a prefill pass on the
chat-templated prompt and a greedy-generation pass capped at $T_{\max}$
tokens. Per-layer pre-hooks registered on the \texttt{mlp} module of every
\texttt{model.layers[l]} capture $\mathbf{h}_t^{(l)}$ at the MLP input,
i.e.\ after the post-attention RMSNorm and before the gate/up projections
(Figure~\ref{fig:transformer_block}).

For every $(l, t)$ pair, the stage computes the seven scalars of
\texttt{LayerEnergyResult}: full energy $E$, entropy $H$, normalised
entropy $\bar{H}$, log-sum-exp term, quadratic term, top score, and
active-row count (rows whose raw score exceeds a fixed threshold
$\tau = 0.1$). It additionally computes the per-token JS and Hellinger
divergences against the prefill last-token reference. The complete set
of per-sample arrays is persisted to a compressed \texttt{\{id\}.npz}
file with the layout listed in Table~\ref{tab:npz_schema}.

\begin{table}[H]
    \centering
    \small
    \begin{tabular}{lll}
        \toprule
        Key & Shape & Description \\
        \midrule
        \texttt{prefill\_energy}        & $[L]$       & Last-token reference energy per layer \\
        \texttt{prefill\_entropy}       & $[L]$       & Last-token reference Shannon entropy per layer \\
        \texttt{prefill\_norm\_entropy} & $[L]$       & As above, normalised by $\log K_l$ \\
        \texttt{prefill\_lse}           & $[L]$       & LSE term of the energy decomposition \\
        \texttt{prefill\_quadratic}     & $[L]$       & Quadratic term $\tfrac{1}{2}\|\mathbf{h}^{(l,\mathrm{ref})}\|^2$ \\
        \texttt{prefill\_top\_act}      & $[L]$       & Max raw score over the bank rows \\
        \texttt{gen\_energy}            & $[L, T_g]$  & Per-layer per-token generation energy \\
        \texttt{gen\_entropy}           & $[L, T_g]$  & Per-layer per-token generation entropy \\
        \texttt{gen\_norm\_entropy}     & $[L, T_g]$  & Normalised generation entropy \\
        \texttt{gen\_lse}               & $[L, T_g]$  & Generation LSE term \\
        \texttt{gen\_top\_act}          & $[L, T_g]$  & Generation max score \\
        \texttt{gen\_js}                & $[L, T_g]$  & JS divergence vs prefill reference \\
        \texttt{gen\_hellinger}         & $[L, T_g]$  & Squared Hellinger distance vs prefill reference \\
        \texttt{token\_ids}             & $[T_g]$     & Generated token ids \\
        \bottomrule
    \end{tabular}
    \caption{Schema of the per-sample \texttt{.npz} artefact produced by
    Stage~2. The companion \texttt{\{id\}.json} records the question, the
    gold answers, the generated text, the generation configuration, and
    the model identifier.}
    \label{tab:npz_schema}
\end{table}

Hidden states themselves are never persisted: the bank multiplication and
the divergences are computed online inside the hook and only the
$[L]$/$[L, T_g]$ scalars survive the forward pass. A diagnostic subset
of $N_{\mathrm{diag}}$ samples additionally writes the raw scores
$\mathbf{M}^{(l)}\mathbf{h}_t^{(l)}$ as a \texttt{\{id\}\_scores.npz}
sidecar, which permits a post-hoc $\beta$ sweep without re-running the
GPU stage.

The stage exposes a sharding interface
(\texttt{shard\_id, num\_shards}) that slices \texttt{samples[k::N]} so
that multiple workers can populate the same trajectory directory in
parallel without lock contention; this is used in the HPC array-job
submission script.

% ----------------------------------------------------------------
\subsection{Stage 3 — Analysis}
\label{subsec:stage3}
% ----------------------------------------------------------------

\textbf{Module: \texttt{pipeline/analysis.py}.} The stage is CPU-only and
reads only the artefacts written by Stage~2. For each sample, it loads the
\texttt{.npz} arrays and the metadata JSON, optionally restricts the
generation arrays to the answer span $\mathcal{T}_q^\star$ (when
\texttt{pool\_over\_answer=True}), and computes the per-layer
$\Delta$-features of Equation~\eqref{eq:per_layer_features} together with
two scalar summaries per sample:
\begin{equation}
    \mathrm{shift}^{(1)}_q
    =
    \tfrac{1}{L}\sum_{l=0}^{L-1}\!\!\left|\Delta E^{(l)}\right|,
    \qquad
    \mathrm{shift}^{(2)}_q
    =
    \!\sqrt{\tfrac{1}{L}\sum_{l=0}^{L-1}\!\!\left(\Delta E^{(l)}\right)^{\!2}}.
    \label{eq:shift_scalars}
\end{equation}
These scalars are not used as detection features but are reported as
interpretable per-sample summaries in the analysis JSON.

The output \texttt{analysis.json} aggregates the per-layer arrays into
$\mathrm{mean}_l$, $\mathrm{std}_l$ across samples, the score summary,
and an optional per-category breakdown (for TruthfulQA's $38$ categories,
which are propagated as \texttt{DataSample.category}).

% ----------------------------------------------------------------
\subsection{Stage 4 — Visualisation}
\label{subsec:stage4}
% ----------------------------------------------------------------

\textbf{Module: \texttt{visualization/visualize\_analysis.py}.} The stage
renders Matplotlib plots from \texttt{analysis.json}: per-layer mean
$\Delta E^{(l)}$, $\Delta \bar{H}^{(l)}$,
$\overline{\mathrm{JS}}^{(l)}$, $\overline{H^2_{\mathrm{Hel}}}{}^{(l)}$
profiles with shaded standard deviations; per-category bar plots; and
the distribution of the per-sample shift scalars. The figures are saved
as PNGs alongside the analysis artefact.

% ----------------------------------------------------------------
\subsection{Stage 5 — Labelling}
\label{subsec:stage5}
% ----------------------------------------------------------------

\textbf{Module: \texttt{labeling/}.} The stage walks every
\texttt{\{id\}.json} produced by Stage~2 and computes both the F1
heuristic label and the LLM-judge label
(Section~\ref{subsec:labels}). The judge backend is configurable:
when \texttt{api\_base} is set, an OpenAI-compatible REST call is made
with the locked template; otherwise, a local HuggingFace model is loaded
for air-gapped HPC use. The judge's verbatim output, the parsed label,
the template hash, and the model identifier are persisted to
\texttt{\{id\}\_labeled.json}. After all samples are labelled,
\texttt{calibration.json} records the inter-rater agreement
(Cohen's $\kappa$) between the heuristic and the judge.

% ----------------------------------------------------------------
\subsection{Stage 6 — Evaluation}
\label{subsec:stage6}
% ----------------------------------------------------------------

\textbf{Module: \texttt{evaluation/}.} The stage joins the per-sample
features written by Stage~3 with the labels written by Stage~5 and
produces:
\begin{enumerate}[leftmargin=*]
    \item per-layer AUROC tables for every feature in
          $\boldsymbol{\phi}_q$ against $z_q$;
    \item the logistic-regression probe of
          Section~\ref{sec:proposed_model} with nested cross-validation;
    \item AUROC tables for the baselines of
          Section~\ref{subsec:baselines};
    \item the post-hoc $\beta$ sweep over the diagnostic subset, when
          raw scores were saved;
    \item a self-contained HTML report with embedded plots.
\end{enumerate}

% ----------------------------------------------------------------
\subsection{Orchestration}
\label{subsec:orchestration}
% ----------------------------------------------------------------

\textbf{CLI: \texttt{hopfield-llm}.} Each stage is exposed as a
sub-command. The orchestrator sub-command \texttt{run-experiment} reads a
single YAML configuration, runs Stages~1–4 with full experiment
tracking (timestamped output directory, git commit hash, configuration
snapshot, timing metadata, structured logs), and writes them inside
\texttt{outputs/\{timestamp\}\_\{experiment\_name\}/}. Stages~5–6 are
explicitly opt-in and run separately, because the labelling stage incurs
external API cost and the evaluation stage benefits from being re-run
multiple times against different label files. The implementation of this
tracking layer lives in \texttt{utils/tracker.py} and is shared across
both local and HPC submissions.

% =============================================================
\section{Testing}
\label{sec:testing}
% =============================================================

Three test tiers protect the pipeline. They are organised as
\texttt{pytest} suites and run by continuous integration on every change.

\paragraph{Unit tests (\texttt{tests/unit/}).} These cover the pure
mathematical primitives: the four terms of the Hopfield energy on
hand-computed inputs; the KL/JS/Hellinger functions on known cases
($p = q$ gives zero; disjoint support saturates Hellinger to one;
$\mathrm{JS} \le \log 2$); the bank-extraction shape assertions for
$\mathbf{W}_{\mathrm{gate}}^{(l)} \odot \mathbf{V}^{(l)}$; and the token-F1
heuristic on hand-labelled toy strings.

\paragraph{Integration tests (\texttt{tests/integration/}).} These run
the four core stages end-to-end on a tiny model and a fixed five-sample
slice of TruthfulQA, asserting that the produced \texttt{analysis.json}
contains the expected keys, that the per-layer arrays have the expected
shapes, and that re-running the same configuration twice produces
bit-identical \texttt{analysis.json} (modulo timestamps).

\paragraph{Pipeline smoke test.} A \texttt{scripts/local/test\_pipeline.sh}
script runs all six stages on the five-sample slice with a public
$\le 1.5$\,B-parameter model, verifies that every expected artefact
exists, and confirms that the chat-template invariance
(Section~\ref{subsec:query_id}) holds: the prefill token sequence and the
prefix consumed by generation must be identical.

The smoke test plus the integration suite are the gate for every commit
that touches files inside \texttt{src/hopfield\_llm/}.

% =============================================================
\section{Proposed Detection Model}
\label{sec:proposed_model}
% =============================================================

% ----------------------------------------------------------------
\subsection{Per-layer logistic probe}
\label{subsec:probe}
% ----------------------------------------------------------------

The proposed detection model is a regularised logistic-regression probe
on the per-layer feature vector $\boldsymbol{\phi}_q \in \mathbb{R}^{4L}$
of Equation~\eqref{eq:feature_vector}. Let $\sigma(\cdot)$ denote the
logistic function. The probe estimates
\begin{equation}
    \hat{z}_q
    =
    \sigma\!\left(\boldsymbol{\phi}_q^{\!\top} \boldsymbol{w} + b\right),
    \qquad
    \boldsymbol{w} \in \mathbb{R}^{4L},\;
    b \in \mathbb{R},
    \label{eq:logistic_probe}
\end{equation}
fit by minimising the L2-regularised binary cross-entropy
\begin{equation}
    \min_{\boldsymbol{w},\,b}\; \tfrac{1}{N}\!\sum_{q=1}^N
    \mathrm{BCE}\!\left(z_q,\,\hat{z}_q\right)
    \;+\;
    \tfrac{1}{2C}\,\|\boldsymbol{w}\|_2^2,
    \label{eq:logistic_objective}
\end{equation}
where $C > 0$ is the inverse-regularisation strength. A standardisation
step ($\mathrm{StandardScaler}$, fit on the training fold only) precedes
the logistic regression. Missing values, which arise when an extraction
warning fires for a single layer, are imputed with the per-column median
of the training fold.

The choice of a linear probe is deliberate: it commits the methodology
to claims that are interpretable per layer (every coefficient in
$\boldsymbol{w}$ maps to a single layer–metric pair), and it places a
strong inductive bias against absorbing surface-feature confounds that a
non-linear classifier could trivially memorise from $4L$ features.

% ----------------------------------------------------------------
\subsection{Held-out evaluation protocol}
\label{subsec:eval_protocol}
% ----------------------------------------------------------------

The dataset of $N$ labelled samples is partitioned, stratified on $z_q$
and with a fixed seed, into a \emph{train}/\emph{validation}/\emph{test}
split of $70/15/15$. The test split is locked before the probe is
fit and is not inspected until the very last step. The hyperparameter
$C$ is chosen from the grid $\{0.01,\,0.1,\,1.0,\,10.0\}$ by inner
$3$-fold cross-validation on the training fold; the outer evaluation is
$5$-fold stratified cross-validation on \emph{train} $\cup$ \emph{val}.
With the best $C$ selected, the probe is finally re-fit on
\emph{train} $\cup$ \emph{val} and evaluated once on the held-out test
split.

This separation is what makes the reported AUROC a genuine generalisation
estimate rather than a cross-validated optimistic statistic.

% ----------------------------------------------------------------
\subsection{Single-feature and per-layer reports}
\label{subsec:single_feature_reports}
% ----------------------------------------------------------------

Alongside the multivariate probe, the methodology reports two
complementary diagnostics. The first is the \emph{per-layer AUROC} of
each individual feature $f \in \{\Delta E, \Delta \bar{H},
\overline{\mathrm{JS}}, \overline{H^2_{\mathrm{Hel}}}\}$ at every layer
$l$. The second is the \emph{single-feature unsupervised AUROC} of the
best-layer score for each $f$, computed without any logistic-regression
fitting. The first diagnostic localises the signal along the depth axis;
the second isolates the unsupervised contribution of the Hopfield-derived
features from the supervised contribution of the probe.

% ----------------------------------------------------------------
\subsection{Baselines}
\label{subsec:baselines}
% ----------------------------------------------------------------

The detection model is compared against a fixed baseline suite that
spans the families normally cited in the hallucination-detection
literature.

\begin{itemize}[leftmargin=*]
    \item \textbf{Token-level uncertainty.} Mean sequence log-probability,
          mean per-token entropy, and maximum per-token entropy of
          $\hat{y}$ under teacher forcing.
          (Module: \texttt{evaluation/baselines.py}.)
    \item \textbf{P(True).} A teacher-forced log-probability that the
          model assigns to the token ``True'' after a follow-up prompt of
          the form ``Is the following answer correct? Answer with a single
          token: True or False.''
          (Same module.)
    \item \textbf{Semantic Entropy} \citep{farquhar2024detecting}.
          $N_s = 10$ stochastic completions per prompt at temperature
          $T_s = 0.5$ are clustered by bidirectional NLI entailment and
          the entropy of the cluster distribution is taken as the score.
          (Module: \texttt{evaluation/semantic\_entropy.py}.)
    \item \textbf{HalluField} \citep{vu2025hallufield}. The logit
          field-theoretic baseline: two teacher-forced forward passes
          with perturbed sampling temperature are used to compute the
          free-energy variation $\delta F_q$.
          (Module: \texttt{evaluation/hallufield.py}.)
    \item \textbf{INTRA} \citep{vazhentsev2026leveraging}. A lightweight
          classifier on layer-pooled residual hidden states. The hidden
          states are read from the same forward pass as the Hopfield
          features, so the cost is fully amortised.
          (Module: \texttt{evaluation/intra.py}.)
\end{itemize}

All baselines are scored under the identical train/val/test split as the
probe. The aim is not to beat each baseline outright; it is to test
whether the Hopfield-derived features add information \emph{on top of}
the token-level uncertainty baselines, which is the contribution that
the central hypothesis claims.

% =============================================================
\section{Analysis Method}
\label{sec:analysis_method}
% =============================================================

This section specifies the statistical procedures used to support every
claim made in Chapter~\ref{chap:experiments}. All procedures are
implemented in \texttt{evaluation/} and run by Stage~6.

% ----------------------------------------------------------------
\subsection{Primary endpoint}
\label{subsec:primary_endpoint}
% ----------------------------------------------------------------

The primary endpoint is the held-out AUROC of the probe of
Equation~\eqref{eq:logistic_probe} on the locked test split, computed
once, with a non-parametric bootstrap $95\%$ confidence interval
estimated from $B = 10\,000$ resamples of the test predictions.

% ----------------------------------------------------------------
\subsection{Per-layer null distribution with scan correction}
\label{subsec:scan_correction}
% ----------------------------------------------------------------

Per-layer AUROC plots are scan statistics: reporting only the maximum
AUROC across $L$ layers inflates the type-I error. The methodology uses
a permutation null in which the label vector is shuffled and the maximum
per-layer AUROC is recomputed; this is repeated $B_{\mathrm{perm}} =
1000$ times to form a null distribution of the scan maximum. A claim of
``layer $l^\star$ shows AUROC above chance'' is supported only if the
observed maximum exceeds the $95$th percentile of this null.

% ----------------------------------------------------------------
\subsection{Within-category permutation null}
\label{subsec:within_category_null}
% ----------------------------------------------------------------

Datasets with categorical structure (TruthfulQA) admit a base-rate
exploit: a detector that simply scores ``Misconceptions'' high will beat
chance if hallucination rates differ across categories. To control for
this, the methodology computes a within-category permutation null in
which labels are permuted only within each category, and reports the
$\Delta\mathrm{AUROC}$ of the probe over this stratified null in
addition to the unconditional one.

% ----------------------------------------------------------------
\subsection{Effect size}
\label{subsec:effect_size}
% ----------------------------------------------------------------

Beyond AUROC, the methodology reports Cliff's $\delta$ between the
predicted-probability distributions of the two classes, with a
bootstrap $95\%$ confidence interval. AUROC measures ranking quality;
$\delta$ measures the standardised separation of the two distributions
and is the effect-size metric appropriate for non-parametric two-sample
comparisons.

% ----------------------------------------------------------------
\subsection{Paired-bootstrap nested-model test}
\label{subsec:paired_bootstrap}
% ----------------------------------------------------------------

The central claim that the Hopfield features add information over a
baseline is tested by the paired-bootstrap nested-model procedure of
\citet{efron1994introduction}. Let $\hat{s}^{\mathrm{full}}_q$ and
$\hat{s}^{\mathrm{base}}_q$ denote the probe scores produced by the
full feature set $\boldsymbol{\phi}_q$ and by the baseline-only feature
subset. For each of $B = 10\,000$ paired bootstrap resamples of the
test set, $\Delta\mathrm{AUROC} = \mathrm{AUROC}^{\mathrm{full}} -
\mathrm{AUROC}^{\mathrm{base}}$ is recomputed; the $95\%$ confidence
interval of this $\Delta$ is reported. A positive lower bound is the
formal condition under which the Hopfield contribution is claimed to be
significant.

% ----------------------------------------------------------------
\subsection{Multiple-comparison correction}
\label{subsec:multiple_comparison}
% ----------------------------------------------------------------

Secondary hypotheses are tested with Bonferroni correction at family
$\alpha = 0.05$. Per-layer AUROC plots are exploratory and are not
subject to Bonferroni; they are protected only by the scan-corrected
null of Section~\ref{subsec:scan_correction}.

% ----------------------------------------------------------------
\subsection{Pre-registration}
\label{subsec:preregistration}
% ----------------------------------------------------------------

Before the final pipeline is executed, the analysis plan — the primary
endpoint, the success threshold, the layer-band of interest predicted
from the theoretical framework (mid-to-late MLP layers, roughly the
$50$–$85\%$ depth band), the baseline suite, and the multiple-comparison
strategy — is committed to the repository as
\texttt{docs/preregistration.md} and tagged with a git release. The
intent of pre-registration is to make every later claim falsifiable: any
deviation from the plan must be documented as a deviation rather than
silently absorbed into the analysis.

% =============================================================
\section{Experimental Configuration}
\label{sec:experimental_config}
% =============================================================

This section specifies the configuration that the final pipeline will be
executed under. The configuration is consistent with the design of
Sections~\ref{sec:solution_design}–\ref{sec:proposed_model}; nothing in
this section overrides earlier design decisions, and no new pipeline
stages are introduced.

% ----------------------------------------------------------------
\subsection{Models}
\label{subsec:models}
% ----------------------------------------------------------------

Two instruction-tuned decoder-only models are evaluated:

\begin{itemize}[leftmargin=*]
    \item \textbf{Qwen2.5-3B-Instruct} \citep{qwen2025qwen25}. $L = 36$
          layers, $D = 2048$, $D_{\mathrm{ff}} = 11008$.
          Primary model for hyperparameter selection and the headline
          result.
    \item \textbf{Llama-3.2-3B-Instruct} \citep{grattafiori2024llama3}.
          $L = 28$ layers, $D = 3072$, $D_{\mathrm{ff}} = 8192$.
          Replication model for cross-architecture transfer.
\end{itemize}

Both models are loaded with $4$-bit NF4 quantisation
(\texttt{bitsandbytes}) on local GPUs and in BF16 on the HPC A100. The
quantisation only affects the forward pass; bank extraction always
dequantises the relevant weight matrices to \texttt{float32} on CPU
before the element-wise product of
Equation~\eqref{eq:gated_key_bank}.

% ----------------------------------------------------------------
\subsection{Datasets}
\label{subsec:datasets}
% ----------------------------------------------------------------

Three short-form factoid QA datasets are used:

\begin{itemize}[leftmargin=*]
    \item \textbf{TruthfulQA} \citep{lin2022truthfulqa}. $817$ adversarial
          questions across $38$ categories. Used as the primary dataset
          because (i) it carries a curated set of correct and incorrect
          reference answers, (ii) it is the de facto benchmark for
          hallucination detection on factoid QA, and (iii) its
          category metadata permits the within-category null of
          Section~\ref{subsec:within_category_null}.
    \item \textbf{TriviaQA} (\texttt{rc.nocontext} split,
          \citealt{joshi2017triviaqa}). Natural-distribution factoid
          questions, used as the cross-dataset transfer evaluation.
          $500$ samples drawn at fixed seed.
    \item \textbf{Natural Questions} (short-answer subset,
          \citealt{kwiatkowski2019natural}). Held in reserve as an
          additional out-of-distribution check if budget allows; not part
          of the primary analysis.
\end{itemize}

For every dataset, sample identifiers are deterministic functions of the
upstream identifier so that artefacts can be traced back to the source
record.

% ----------------------------------------------------------------
\subsection{Bank, probe, and generation hyperparameters}
\label{subsec:hyperparameters}
% ----------------------------------------------------------------

Table~\ref{tab:hyperparameters} lists every hyperparameter of the
final pipeline. These values are fixed before execution and are not
varied across the runs that produce the headline result.

\begin{table}[H]
    \centering
    \small
    \begin{tabular}{p{0.36\linewidth} p{0.28\linewidth} p{0.28\linewidth}}
        \toprule
        Parameter & Value & Provenance \\
        \midrule
        \multicolumn{3}{l}{\emph{Memory bank}} \\
        \hspace{1em}\texttt{bank}                & \texttt{gated\_key}        & Eq.~\eqref{eq:gated_key_bank} \\
        \hspace{1em}\texttt{normalize}           & \texttt{True} (row L2)     & §\ref{subsec:memory_def} \\
        \hspace{1em}\texttt{hook\_target}        & \texttt{mlp\_input}        & §\ref{subsec:query_id} \\
        \multicolumn{3}{l}{\emph{Energy and divergences}} \\
        \hspace{1em}\texttt{energy\_mode}        & \texttt{dot}               & §\ref{subsec:retrieval_scoring} \\
        \hspace{1em}\texttt{normalize\_query}    & \texttt{False}             & §\ref{subsec:retrieval_scoring} \\
        \hspace{1em}\texttt{threshold} (active-row count) & $0.1$              & — \\
        \hspace{1em}\texttt{beta} ($\beta^\star$)& calibrated, target $\bar{H}^\star = 0.55$ & §\ref{subsec:beta_calibration} \\
        \hspace{1em}\texttt{calibration\_samples}& $30$                       & §\ref{subsec:beta_calibration} \\
        \multicolumn{3}{l}{\emph{Generation}} \\
        \hspace{1em}Decoding                     & greedy                     & §\ref{subsec:constraints} \\
        \hspace{1em}\texttt{max\_new\_tokens}    & $50$                       & — \\
        \hspace{1em}Chat template                & model default              & §\ref{subsec:query_id} \\
        \multicolumn{3}{l}{\emph{Analysis pooling}} \\
        \hspace{1em}\texttt{pool\_over\_answer}  & \texttt{True}              & §\ref{subsec:feature_design} \\
        \multicolumn{3}{l}{\emph{Probe}} \\
        \hspace{1em}Regulariser                  & L2                          & Eq.~\eqref{eq:logistic_objective} \\
        \hspace{1em}$C$ grid                     & $\{0.01, 0.1, 1, 10\}$     & §\ref{subsec:probe} \\
        \hspace{1em}Inner CV                     & $3$-fold stratified         & §\ref{subsec:eval_protocol} \\
        \hspace{1em}Outer CV                     & $5$-fold stratified         & §\ref{subsec:eval_protocol} \\
        \hspace{1em}Split                        & $70/15/15$ (train/val/test) & §\ref{subsec:eval_protocol} \\
        \multicolumn{3}{l}{\emph{Statistics}} \\
        \hspace{1em}Bootstrap resamples          & $10\,000$                   & §\ref{subsec:primary_endpoint} \\
        \hspace{1em}Permutation resamples        & $1000$                      & §\ref{subsec:scan_correction} \\
        \hspace{1em}Multiple-comparison family   & Bonferroni, $\alpha=0.05$  & §\ref{subsec:multiple_comparison} \\
        \multicolumn{3}{l}{\emph{Labelling}} \\
        \hspace{1em}F1 threshold $\tau_{\mathrm{F1}}$ & $0.3$                  & §\ref{subsec:labels} \\
        \hspace{1em}Judge model                  & \texttt{gpt-4o-mini}        & §\ref{subsec:labels} \\
        \hspace{1em}Prompt template              & \texttt{convcoa\_v1.txt}    & §\ref{subsec:labels} \\
        \multicolumn{3}{l}{\emph{Determinism}} \\
        \hspace{1em}Random seed                  & $42$                       & — \\
        \hspace{1em}Generation                   & deterministic (greedy)      & §\ref{subsec:constraints} \\
        \hspace{1em}Bank normalisation seed      & not applicable             & — \\
        \bottomrule
    \end{tabular}
    \caption{Final pipeline hyperparameters. The provenance column points
    at the section in which each choice is justified.}
    \label{tab:hyperparameters}
\end{table}

% ----------------------------------------------------------------
\subsection{Storage budget}
\label{subsec:storage}
% ----------------------------------------------------------------

For Qwen2.5-3B on $500$ TruthfulQA samples ($L=36$, $K=11008$,
$T_g \le 50$), each \texttt{\{id\}.npz} occupies $30$–$60$\,KB after
compression, the optional \texttt{\{id\}\_scores.npz} sidecar weighs
$\approx 38$\,MB per sample when full scores are kept, and the
\texttt{banks.pt} artefact is $\approx 3$\,GB. The default
configuration writes full scores only for the first $50$ samples
(\texttt{scores\_subset\_size = 50}), bounding the total trajectory
directory to $\approx 2$\,GB. Hidden states are never persisted; they
live only in GPU memory during the forward pass.

% ----------------------------------------------------------------
\subsection{Final execution plan}
\label{subsec:execution_plan}
% ----------------------------------------------------------------

The final pipeline is executed as a single, ordered sequence. Each step
maps directly to a CLI invocation that consumes the artefacts of the
previous step. The sequence is reproducible: given the same model
checkpoint, the same dataset split, and the same seed, it produces
bit-identical artefacts up to floating-point noise from the BF16/4-bit
quantisation paths.

\begin{enumerate}[leftmargin=*]
    \item \textbf{Build the gated-key bank for Qwen2.5-3B.}
        \texttt{hopfield-llm build-banks --model qwen25\_3b --bank
        gated\_key --output outputs/qwen25\_3b/banks.pt}. Produces
        \texttt{banks.pt} and \texttt{banks\_metadata.json}.
    \item \textbf{Calibrate $\beta^\star$ on a 30-question warm-up.}
        Performed inline by Stage~2 when
        \texttt{calibrate\_beta=true}; writes
        \texttt{calibration\_beta.json} into the run directory.
    \item \textbf{Capture trajectories on TruthfulQA.}
        \texttt{hopfield-llm run-trajectory --model qwen25\_3b --dataset
        truthfulqa --banks outputs/qwen25\_3b/banks.pt --output
        outputs/qwen25\_3b/truthfulqa/trajectories}. Writes one
        \texttt{.npz}/\texttt{.json} pair per sample.
    \item \textbf{Analyse trajectories with answer-span pooling.}
        \texttt{hopfield-llm analyze --trajectories \dots
        --pool-over-answer --output analysis.json}.
    \item \textbf{Render the Stage~4 plots.}
        \texttt{hopfield-llm visualize --analysis analysis.json
        --output plots/}.
    \item \textbf{Label with the F1 heuristic and the LLM judge.}
        \texttt{hopfield-llm label --trajectories \dots --output
        labels/ --judge-model gpt-4o-mini}. Writes per-sample labels
        and \texttt{calibration.json}.
    \item \textbf{Run the baselines.}
        \texttt{python -m hopfield\_llm.evaluation.baselines}
        (token-uncertainty + P(True)) on the same trajectories;
        \texttt{python -m hopfield\_llm.evaluation.semantic\_entropy}
        for the sampling baseline; \texttt{python -m
        hopfield\_llm.evaluation.hallufield} for HalluField;
        \texttt{python -m hopfield\_llm.evaluation.intra} for INTRA.
        All four write per-sample sidecars into the trajectory
        directory.
    \item \textbf{Evaluate.}
        \texttt{hopfield-llm evaluate --trajectories \dots --labels
        labels/ --output report.html}. Performs the
        train/validation/test split, fits the probe with nested CV,
        computes per-layer AUROCs, paired-bootstrap nested-model tests,
        bootstrap CIs, permutation nulls, and writes the HTML report.
    \item \textbf{Replicate on Llama-3.2-3B.} Steps 1–8 are repeated
        with \texttt{--model llama32\_3b}. The fitted probe from the
        Qwen run is additionally applied zero-shot to the Llama
        features as the cross-model transfer check.
    \item \textbf{Cross-dataset transfer.} Steps 3–8 are repeated for
        each model with \texttt{--dataset triviaqa}. The Qwen probe
        fitted on the TruthfulQA train+val split is evaluated on the
        TriviaQA test split without retraining.
\end{enumerate}

The ordering of the steps reflects the dependency graph among the
artefacts: bank $\rightarrow$ trajectories $\rightarrow$ analysis,
trajectories $\rightarrow$ labels, and \{analysis, labels, baselines\}
$\rightarrow$ evaluation. Re-running an earlier step does not require
re-running upstream steps; re-running a later step does not invalidate
earlier artefacts.

% ----------------------------------------------------------------
\subsection{Reproducibility artefacts}
\label{subsec:reproducibility}
% ----------------------------------------------------------------

Each run directory contains:

\begin{itemize}[leftmargin=*]
    \item the resolved YAML configuration, snapshotted at execution time;
    \item the git commit hash and a flag indicating whether the working
          tree was clean;
    \item the environment fingerprint (Python version, \texttt{torch}
          version, \texttt{transformers} version, \texttt{bitsandbytes}
          version, CUDA version);
    \item the per-stage wall-clock timing recorded by
          \texttt{ExperimentTracker};
    \item the chat-template string verbatim, including any model-specific
          control tokens;
    \item the random seed and the deterministic ordering of the dataset
          slice consumed.
\end{itemize}

These artefacts are sufficient to re-execute the methodology on a
different host and to verify the numerical reproducibility claim.

% =============================================================
\section{Summary}
\label{sec:method_summary}
% =============================================================

The methodology turns the theoretical framework of
Chapter~\ref{chap:theory} into a concrete, auditable pipeline. The probe
treats the hidden state entering each SwiGLU sublayer as a query against
a layer-specific bank whose rows are the element-wise product
$\mathbf{W}_{\mathrm{gate}}^{(l)} \odot \mathbf{V}^{(l)}$ of the gate and
value projections, computes the Modern Hopfield retrieval energy and
divergences against the prefill last-token reference, pools the resulting
per-token quantities over the answer span, and feeds the per-layer
$\Delta$-features into an L2-regularised logistic regression evaluated on
a locked $15\%$ test split. The pipeline writes typed artefacts at every
stage, is sharded for HPC use, and is supported by unit, integration, and
end-to-end smoke tests. The final configuration of
Section~\ref{sec:experimental_config} is the one that
Chapter~\ref{chap:experiments} will execute and report against
pre-registered baselines and pre-registered statistical procedures.

% =============================================================
% END OF CHAPTER
% =============================================================
